\documentclass[11pt]{article}
\usepackage{amsmath,amssymb,amsthm,geometry}
\geometry{margin=.9in}\setlength{\emergencystretch}{1.3em}
\newtheorem{theorem}{Theorem}\newtheorem{lemma}{Lemma}
\newcommand{\dd}{\,d}
\title{A trained critical bridge at first strong half mass}
\author{AN04 proof increment}\date{Version 1}
\begin{document}\maketitle

\paragraph{Clock and signed observable.}
Work in P's normalized time with aligned balanced isotropic continuum
initialization, $S_0=cq^{1/\omega}$, fixed $c>0$,
$13/20\le\rho\le7/10$ fixed, $\lambda=\rho^2$, $\omega=(1-\lambda)/2$.
The data are the equally weighted Gaussians
$P_1=N(e_1,q^2I)$ and $P_2=N(\rho e_2,q^2I)$.
Write $\alpha$ for the initial angle with law $\dd l=\dd\alpha/(2\pi)$,
$U=\sqrt m\,a_\theta$, $W=\sqrt m\,b_\psi$ for the two vectors,
and $M_q=\int_{-\pi/2}^{\pi/2}m\dd l$ for the radial mass of fixed
right-half ancestry. Use the scaled current angles $x=\theta/q$,
$y=\psi/q$ on the right-facing chart. Let $a=\rho K(\rho a)$,
$K(t)=\phi(t)+t\Phi(t)$ and fix $\zeta>a$.
In this chart the probe is active for $x<\zeta$ at bounded coordinates.
For a fixed label band $B$, put
$f_B(D)=\int_B W(U^TD)_+\dd l$ and
$C_B=\mathbb E[f_B(D)D^T{\bf1}_{a_\theta^TD>0}]$.
Its actual fitted contribution to the two scaled receiver velocities is
\[
  (v_x^B,v_y^B)=-q^{-1}
       (b_\psi^TC_Ba_\theta^\perp,(b_\psi^\perp)^TC_Ba_\theta).
\]
All expectations include the mixture weights. These are contributions
at the actual trained state, with all donors present during training.

\begin{theorem}[Trained flow]\label{an04hbdv1:main}
For every nonempty compact interior Q4 interval $C$ of positive length,
there exist a fixed compact interior Q1 interval $B$ of positive
length and $\gamma>0$ such that for all sufficiently small $q$:
\begin{enumerate}
\item The first time $t_h$ with $M_q=1/2$ exists and
$t_h=\log(4/S_0)+o(1)$.
\item At $t_h$ every receiver in $C$ is strictly active and every
donor in $B$ is strictly inactive, with fixed scaled gate margins.
The radial mass of $B$ tends to
$m_B^*=2\int_B\cos^2\alpha\dd l>0$.
\item At each receiver in $C$, $v_x^B\le-\gamma$ and
$v_y^B\le-\gamma$ at $t_h$. The same pair of strict signs holds for
the difference between these velocities and those produced by an
equal mass placed at the receiver's current coordinates.
Compared with concentrating \emph{all} right-half mass at that receiver,
the full trained angular field also has both components smaller by at
least $\gamma$ in scaled units. These three velocity comparisons hold
throughout $[t_h-1/4,t_h]$.
\end{enumerate}
\end{theorem}

\paragraph{Models and dependencies.}
P exactflow, representation, earlybounds and targetcomparison apply
to this Gaussian initialization; the early angular and radial comparison
is used only below its small envelope cap. CFS laws and RPS q4 are
target-only, positive-noise sector limits at $T_A=\log(A/S_0)$ for
fixed $A,c>0$, uniform on compact interior labels. CTE early and profile
are trained compact-Q1 estimates on the small cap. The compact-chart
Gaussian target and per-donor formulas below are the individual kernels
in MOMC donorkernels, whose expansions require a bounded receiver and
bounded donor chart, not ten-power initialization. We treat all other
donors separately. WIT popsystem and RPS rank specify the limiting
field; the passage to this field and its ingoing conditions are proved
here. The time interval below ends at $S_0 e^t=8$, with no weak-switch
or weak-seed assumption.

\section{Physical bounds before taking chart limits}
Use $\tau=t+\log S_0$, $A=e^\tau$, $H=\sqrt{S_0}e^{a_*t}$,
$a_*=1/2+q^2$, $d_*=\lambda/2+q^2$, $E=e^{-\omega t}+q$.
Through $\tau\le\log8$ one has $H^2=Ae^{2q^2t}\le9$ for small $q$.
P's spectral envelope gives $S=\int m\dd l\le H^2$ throughout.
The generic bound $|(\log m)'|\le 2a_*(1+S)$ follows from exactflow
and $\|C_f\|\le a_*S$. Integration of the upper bound yields
$\sqrt m\le H\exp(H^2/2)\le C H$ for every label. Here and below
$C$ can change but is independent of $q$ and label. We claim
\begin{equation}\label{an04hbdv1:physical}
 B_2:=\sup_\alpha\max(|U_2|,|W_2|)\le C H E .
\end{equation}
Indeed $U'=R^TW$, $W'=RU$, $R=A_\theta-C_f$, with
$|(A_\theta)_{12}|\le\mathbb E|D_1D_2|\le Cq$,
$(A_\theta)_{22}\le d_*$. The fitted output obeys
$|f_1|\le CH(H|D_1|+B_2|D_2|)$ and $|f_2|\le CHB_2|D|$,
so $|(C_f)_{12}|+|(C_f)_{21}|\le C(HB_2+qH^2)$.
Hence $D^+B_2\le(d_*+CH^2)B_2+CqH(1+H^2)$.
The integral of $H^2$ is at most $9/(2a_*)$, $B_2(0)\le\sqrt{S_0}$,
and $a_*-d_*=\omega>0$. Scalar integration proves
\eqref{an04hbdv1:physical}. This extends the physical bound beyond
the small-cap range of CTE without extending its angular comparison.

At any fixed small nominal cap $A_0$, the entire left ancestry has
mass per label $O_{A_0}(q^2(1+t)^2+S_0)$.
For clarity, in target-only dynamics the component equation of
PROFILE coordinate gives for positive parts
$p_j'\le(\mu_j^2/2+q^2)p_j+CqH$, $\mu=(1,\rho)$, and negative
parts grow at most at rate $q^2$. With $p_1(0)=0$ on the left,
$p_1\le CqtH$, $p_2\le\sqrt{S_0}e^{d_*t}+CqH$ and the negative
parts are at most $\sqrt{S_0}e^{q^2t}$. Square and use
$e^{-2\omega t}=q^2(c/A_0)^{2\omega}$ at observation.
P targetcomparison multiplies the weight by at most $e^{4H^2}$
on this small cap. The generic logarithmic bound propagates the
estimate from $A_0$ to $8$ over a fixed elapsed interval. In
particular the entire left contribution divided by $q$ tends to
zero there, even using just its mass and $\|C_{\rm left}\|\le
a_*\int_{\rm left}m\dd l$.

On a compact interior right label set and a compact $\tau$ interval
the radial weights are bounded below by a positive constant. Start at
a smaller fixed cap with CFS laws or RPS q4 and P targetcomparison,
then integrate the generic lower logarithmic bound over finite time.
At that starting cap both angles are on the right chart: CTE supplies
this on Q1 and the Q4 version is justified below. Equation
\eqref{an04hbdv1:physical} then keeps both angles on that chart
($\sqrt m$ is bounded below and the sines are $O(q)$), with
bounded scaled coordinates. This assertion is only label-compact.

\paragraph{Q4 entry.}
The proof of CTE early for an individual receiver requires just
$|\sin\widehat\theta|\le C_I E$ on the target orbit in addition to
the all-sector physical bound and early alignment guard. For compact
Q4, while $\widehat\theta<0$, PROFILE angle gives
$\widehat\theta'\ge-\sin\widehat\theta\cos\widehat\theta/2$,
hence $\tan|\widehat\theta|\le e^{-t/2}\tan|\alpha|$.
At zero the field points up, and at a sufficiently large positive
angle $qL$ it points down, by its fixed-chart limit
$[-L+\rho K(\rho L)]/2<0$. This proves the sine bound from
initialization. To spell out the comparison, on the small cap the
antisymmetric fitted projection at alignment is $O(H^2E)$ by
\eqref{an04hbdv1:physical}; away from alignment add
$CH^2|\psi-\theta|$. The target alignment term is inward under
P's lift guard. Integration gives $|\psi-\theta|\le CH^2E$.
For $e=\theta-\widehat\theta$, $|v_0'|\le a_*$ and the aligned
fitted projection is $O(H^2(E+|\sin\theta|))$, giving
$D^+|e|\le(a_*+CH^2)|e|+C_IH^2E$. Integration from zero uses
$2a_*-\omega-a_*=d_*>0$ and gives the same $C_IH^2E$ error.
Thus, by RPS q4, at each fixed small $A$ both Q4 scaled angles
are within $C_I A[1+(c/A)^\omega]+o_{q;A,I}(1)$ of $a$.
On Q1 the same bound about $H_{\rm CFS}^{-1}((c/A)^\omega\tan\alpha)$
is CTE profile. These estimates concern each compact interior set;
their constants need not stay bounded at a vertical initial axis.

\section{All donors in the limiting equations}
The following argument uses subsequences to avoid assuming a unique
full profile arriving from infinite negative $\tau$. All the inequalities
used at half mass will hold for every subsequential profile.

\paragraph{Compact kernels and endpoint costs.}
For a bounded-chart receiver and bounded-chart donor at
$(\xi,\eta)$ the limiting per-unit-mass fitted contributions,
with the sign of the fitted residual already included, are
\[
\begin{array}{c|cc}
 &2v_x&2v_y\\ \hline
 P_1\text{ term}&x&y-\eta\\
 P_2\text{ term}&-\rho F(\rho x,\rho\xi)&0
\end{array}
\qquad
 F(t,r)=\phi(\min(t,r))+r\Phi(\min(t,r)).
\]
Here the two rows sum to the doubled mixture velocity.
On $P_1$ the right-facing gates tend to one, while
$a^\perp\cdot D=q(Z_2-x)+O_{L^2}(q^2)$.
On $P_2$, $D=(qZ_1,\rho+qZ_2)$ and the donor activation divided by
$q$ tends to $(Z_1+\rho\xi)_+$. The receiver gate selects
$Z_1>-\rho x$. These limits give the table by Gaussian integration
and are uniform on fixed compact charts (also directly MOMC
donorkernels). The doubled target velocities tend to
$(-x+\rho\phi(\rho x)+\lambda y\Phi(\rho x),
-y+\rho K(\rho x))$; the log radial target and fitted parts tend to
$1$ and $-1$ per unit fitted mass respectively, with the single
target counted once. These follow from the same Gaussian moments.

For any set of right labels of initial length $\eta$, on any compact
$\tau$ interval its physical output satisfies
$|f_1|\le C\eta(|D_1|+q|D_2|)$ and $|f_2|\le C\eta q|D|$ by
\eqref{an04hbdv1:physical}. At a bounded receiver its fitted angular
contribution after division by $q$ is at most $C_R\eta$:
on $P_1$, $a^\perp\cdot D=O_{L^2}(q)$ and
$(b^\perp)^Tf=O_{L^2}(\eta q)$; on $P_2$, $f_1=O_{L^2}(\eta q)$,
$f_2=O_{L^2}(\eta q)$. Dropping gates in absolute bounds is valid.
Radial costs are $O(\eta)$. One may therefore pass first on compact
label interiors, then send the omitted strip lengths to zero. This
accounts for donors potentially much further out than the scaled
receiver chart. The common velocity and weight bounds also give
equicontinuous time paths on each compact interior and compact
$\tau$ interval.

\paragraph{Extraction and entrance.}
For any sequence $q\downarrow0$, push the fixed right Lebesgue label
measure forward to the joint law of label and its three time paths
$(x,y,m)$; keep its unnormalized total label mass $1/2$.
On compact time intervals bounded charts, weight bounds and
equicontinuity on each interior label compact give tightness by
Arzel\textendash Ascoli, discarding arbitrarily small endpoint and axis
strips. Diagonal extraction on integer time windows gives a limit
law of paths on $(-\infty,\log8]$ (one may extend the paths
constantly before initialization for each $q$). The marginal in
label is still $\dd l$. Any finite list of tagged paths, even with
labels varying in a fixed interior compact, can be extracted on the
same windows. Integrals on interior compacts pass by weak convergence
of the joint law, evaluation of continuous paths, and the compact
kernel bounds just proved; use label cutoffs with null boundary.
Uniform time bounds allow integration in time before passing to
the limit. Equivalently apply this argument in the integral equations
with a tagged receiver and a second, independently sampled donor
from the same law. For almost every limiting path the integral
equations hold: the mean absolute residual of the integral equations
on each receiver compact is zero in the limit by compact-kernel
convergence and strip domination (test first at rational endpoints).
They hold also for every extra tagged path.
In particular if $\mu_\tau$ is the law weighted by the current $m$,
$M=\int\dd\mu_\tau$, $Y=\int y\dd\mu_\tau$,
$K_s=\int K(\rho x)\dd\mu_\tau$, and
$S(x)=\int F(\rho x,\rho\tilde x)\dd\mu_\tau$, then
\begin{align}
 (\log m)'&=1-M,\qquad M'=M(1-M),\nonumber\\
 2x'&=-(1-M)x+\rho\phi(\rho x)-\rho S(x)+\lambda y\Phi(\rho x),
 \qquad
 2y'=-(1-M)y-Y+\rho K(\rho x). \label{an04hbdv1:field}
\end{align}
Primes now denote $\tau$ derivatives. The raw moments needed here
are uniformly locally integrable. More precisely
\begin{equation}\label{an04hbdv1:budget}
 |Y|+\rho K_s\le C e^\tau(1+c^\omega e^{-\omega\tau}),
 \qquad \tau\le\log8 .
\end{equation}
Indeed on each right compact, in the limit $m|y|$ and $m|x|$
are bounded by the limits superior of
$2\sqrt m B_2/q\le C H^2 E/q$; the angles are small there.
The same constant works on all compacts since the physical bound
is global. Exhaust the interiors and use $K(r)\le\phi(0)+r_+$.
This is the budget of the actual limit law, including the outliers.

At each early fixed $A$, the target weight laws and
$|\log(m/\widehat m)|\le4H^2$ give
$m=A\cos^2\alpha\, e^{O(A)}$ in the limit, on every label interior.
The constant in the relative bound is global. Thus as $\tau\to-\infty$,
$M/e^\tau\to1/4$ and $m/e^\tau\to\cos^2\alpha$.
These assertions, and the angular entry constraints below, may first
be imposed at a countable dense set of nominal levels, and extended
by continuity on compact time intervals. Integration of the radial
equations fixes the mass and the normalized label weights:
\begin{equation}\label{an04hbdv1:clock}
 M=\frac{e^\tau}{4+e^\tau},\qquad m=4M\cos^2\alpha .
\end{equation}
The same formula holds on tagged paths with the limiting label.
Every subsequence gives the same mass on compact times, uniformly
there by equicontinuity of $M_q$. Before any fixed small $A_0$
the spectral envelope rules out a half hit. The strictly increasing
limit crosses $1/2$ at $\tau_*=\log4$, below $\log8$.
Bracketing on either side and using continuity proves the asserted
first-hit clock. Uniform local time compactness allows every
subsequent evaluation at the moving first hit.

For Q4 the limiting entrance estimate above gives $x,y\to a$
in the remote past. For Q1 it gives
\begin{equation}\label{an04hbdv1:incoming}
 e^{\omega\tau}x,\ e^{\omega\tau}y\ \longrightarrow\
 c^\omega\tan\alpha\qquad(\tau\to-\infty),
\end{equation}
using $H_{\rm CFS}(r)/r\to1$. Both assertions hold uniformly in the
sense of the bounds on any one fixed interior label compact, hence
for its tagged paths as well. No assertion of pointwise input order
in the finite-noise original flow has been used.

\section{The active receiver and the outer band}
\paragraph{Same-field order and the lower atom.}
In \eqref{an04hbdv1:field} both off-diagonal receiver derivatives
are $\lambda\Phi(\rho x)/2>0$ with the donors held fixed. Differences
in this same time-dependent field therefore preserve coordinate order
by the linear difference equation with segment-averaged Jacobian
and nonnegative off-diagonals. The two diagonal entries of the
\emph{doubled} field are
\[
 -(1-M)+\lambda\rho(y-x)\phi(\rho x)
 -\rho^2\phi(\rho x)\int(\rho\tilde x-\rho x)_+\dd\mu_\tau,
 \qquad -(1-M).
\]
Differentiation uses \eqref{an04hbdv1:budget}.
Any two limiting Q4 paths coalesce. Indeed both approach $(a,a)$
as $\tau\to-\infty$. On their connecting segments for all sufficiently
negative $\tau$, the row sums of the doubled Jacobian are at most
$-(1-\lambda)/2$, since $M\to0$, $y-x\to0$, and the integral enters
with a minus sign. The infinity-norm difference therefore decays
forward there at rate at least $(1-\lambda)/4$ by the scalar
upper derivative inequality. Integrate from arbitrarily early bounded
data, then use same-field uniqueness forward. Denote this common
path by $(x_a,y_a)$. Every Q1 path is above it in both coordinates:
\eqref{an04hbdv1:incoming} gives strict separation far enough back
for each such path, and same-field order propagates it. This applies
also to the separately tagged receiver or donor. The exceptional
initial axes have zero label measure.

Both lower-atom coordinates satisfy $x_a,y_a\le a$ through $\log8$.
For $d=a+\eta$ with $\eta>0$, the closed upper region $x_a,y_a\le d$
has strict remote entry. Since all inputs exceed or equal $x_a$,
$S(x_a)=M\phi(\rho x_a)+\rho\Phi(\rho x_a)\int\tilde x\dd\mu_\tau$.
On $x_a=d$, $y_a\le d$, use $\int\tilde x\dd\mu_\tau\ge Md$
in \eqref{an04hbdv1:field}. The doubled velocity is at most
$(1-M)[\rho K(\rho d)-d]<0$. On $y_a=d$, $x_a\le d$, use
$Y\ge Md$ to obtain $2y_a'\le\rho K(\rho d)-d<0$.
The scalar $d-\rho K(\rho d)$ increases strictly and vanishes at
$a$. Both faces, including their corner, are inward. Taking
$\eta\downarrow0$ proves the claim. In particular $x_a<\zeta$.
For a fixed Q4 compact $C$ there is also a lower bound $x_a,y_a\ge-L_C$
near $\tau_*$ common to all subsequences, taking for example $\log2\le\tau\le\log6$
as the neighborhood. It follows directly from
the radial lower bound and \eqref{an04hbdv1:physical} on $C$,
and does not use an order comparison between different fields.

\paragraph{Uniform incoming lower estimate for high Q1 labels.}
Set $w=\min(x,y)$ on a limiting Q1 path. While $w\ge0$,
\[
 2w'\ge -(1-\lambda)w-J(\tau),\qquad
 J(\tau)=C e^\tau(1+c^\omega e^{-\omega\tau}).
\]
Indeed at the smaller input, $\rho\phi(\rho x)+\lambda y\Phi(\rho x)
\ge\rho K(\rho x)$, and $S(x)\le K_s$; at the smaller output
$\rho K(\rho x)\ge\rho K(\rho y)$. Use $\rho K(\rho w)\ge\lambda w$,
$Mw\ge0$ and \eqref{an04hbdv1:budget} on both faces. The minimum is
absolutely continuous and the inequality holds at ties in the
one-sided comparison sense. Let
\[
 J_*=\tfrac12\int_{-\infty}^{\log8} e^{\omega r}J(r)\dd r <\infty.
\]
For any Q1 label with $c^\omega\tan\alpha>J_*$,
integration from the remote past gives
\begin{equation}\label{an04hbdv1:outer}
 w(\tau)\ge e^{-\omega\tau}(c^\omega\tan\alpha-J_*)>0
 \qquad(\tau\le\log8).
\end{equation}
A first zero is excluded by the same inequality up to contact.
This proof takes the incoming limit for each fixed interior label
before selecting the high band. It never requires a uniform
target-only estimate out to the vertical axis. The budget constant
does not depend on the selected band.

\paragraph{Signs and constants at the measured clock.}
Take $p_C=\Phi(-\rho L_C)>0$. Choose a finite $L$ with
\[
 L>\max\{\zeta+2,\ a+2,\ (a+2)/(\lambda p_C)\}
\]
and a positive-length compact Q1 interval $B$ on which
$c^\omega\tan\alpha>J_*+8^\omega L$.
Such an interval exists for every fixed $c>0$. This choice can be
made in the early chart too: \eqref{an04hbdv1:incoming} and the
early estimates, on the now fixed $B$, place both coordinates beyond
the gate at sufficiently small fixed nominal mass.
On $B$, \eqref{an04hbdv1:outer} gives $x,y>L$ at $\tau_*$.
For actual labels in $B$ the same lower bounds with a small margin
loss hold at the first hit and on each fixed prior compact clock
interval: otherwise extract the violating sequence as a tagged path.
The same argument on $C$ gives $x\le a+o(1)$, $y\le a+o(1)$ at
the hit, uniformly; take for example half the positive gap
$\zeta-a$ as the activity margin. Equation~\eqref{an04hbdv1:clock}
gives $m_B^*$ and, by boundedness on $B$, positive lower bounds
for both band moments, with limit inferiors at least $L m_B^*$.

In every subsequential limit at $\tau_*$, write $x=x_a,y=y_a$ and
$(\xi,\eta)$ for the band donor. The exact trained fitted contributions
at the hit converge along that subsequence to the compact Gaussian
kernel integrals
\begin{align}
 2v_x^{B,\lim}
   &=\int_B[x-\rho\phi(\rho x)-\lambda\xi\Phi(\rho x)]\dd\mu,
 &2v_y^{B,\lim}&=\int_B(y-\eta)\dd\mu .\label{an04hbdv1:sign}
\end{align}
Here $B$ denotes the initial labels, gated by their current inputs:
$\xi>x$ permits the exact overlap simplification. The first integrand
is at most $a-\lambda p_C L<-2$, and the second at most $a-L<-2$.
Thus for sufficiently small $q$ both scaled velocities are at most
$-m_B^*/2$ uniformly on $C$. In the equal-mass comparison the
doubled differences instead are
$-\lambda\Phi(\rho x)\int_B(\xi-x)\dd\mu$ and
$-\int_B(\eta-y)\dd\mu$. Both are bounded above by
$-m_B^*(L-a)\min(1,\lambda p_C)$.
For example a common trained margin in the theorem is
\[
 \gamma=\min\{m_B^*/2,\,
          m_B^*(L-a)\min(1,\lambda p_C)/4\}>0 .
\]
All margins follow for the full small-$q$ family by contradiction
and tagged extraction. The proof has not assumed convergence to a
unique critical angular profile.

\paragraph{Accounting for the other donors.}
In \eqref{an04hbdv1:field} the full velocity includes the target
and the fitted field of $B$ \emph{and its complement}. Equivalently,
for a receiver in $C$ at the hit its doubled limiting velocity
is the target vector specified in the kernel calculation plus the
integral of $(x-\rho F(\rho x,\rho\xi),y-\eta)$ over the whole
right label law with current weights. Splitting that integral by
initial labels gives \eqref{an04hbdv1:sign} and an explicit complement
integral. The left error is $o(1)$ in scaled velocity, and the
endpoint-strip bound is $O_C(\eta_0)$ for strip length $\eta_0$,
sent to zero after taking the local chart limits. The raw complementary integral can be order one.
For a second comparison, replace every right donor by a copy of the
receiver, preserving the full right mass and the left population.
The target cancels. On the bottom path all donors have
$\xi\ge x,\eta\ge y$; each doubled difference from the copy equals
$(-\lambda\Phi(\rho x)(\xi-x),-(\eta-y))$ per unit mass.
Thus every right donor outside $B$ adds a nonpositive difference in
both components. The margin from $B$ also holds for the full field
relative to this mass-matched concentrated state. This conclusion
uses limiting order only; finite-$q$ ordering is unnecessary.

In fact the entire tagged compactness argument applies at arbitrary
moving times in $\log2\le\tau\le\log6$. In that range
$M\ge1/3$, so the band has limiting mass at least $2m_B^*/3$,
and \eqref{an04hbdv1:outer} still gives $x,y>L$ on the band.
The strict margins above leave room for precisely the same $\gamma$:
the raw bound is strictly more negative than $-2m_B^*/3$,
and the equal-mass bound is at most
$-m_B^*(L-a)\min(1,\lambda p_C)/3$.
Since $[t_h-1/4,t_h]$ lies inside this range in the shifted
clock for small $q$, integration gives a contribution at most
$-\gamma/4$ in each scaled coordinate on that interval, along any
one of the actual tagged core trajectories. The equal-mass integrals
are evaluated on this same receiver path; converting them to a
difference of evolved probe errors requires solving the response. The moving-clock error in these statements is $o(1)$
by compact time convergence. No power rate or universal numerical
lower bound on the mass of $B$ is asserted.

\paragraph{Handoff.}
This is a trained result at the first strong half-mass time. The
analyzed delay in $\tau$ stays bounded above by $\log8$; the proof
does not continue to a weak-learning time diverging in this clock.
The remaining depth inequality requires the time integral of the
receiver's response to the donor field through weak learning,
with the actual weak seed, current gates, and appropriate reference
state. Two negative instantaneous fitted velocities at half mass
do not provide that integral or the optimized trained minimum.
\end{document}
